% GENERATED FILE. Edit canonical lessons, then run npm run generate:books.
% canonical-source-hash: fnv1a64:3c2cf140eae61510
% canonical-lessons: RU-C244-drevniy, RU-C244-gromkiy, RU-C244-yarkiy, RU-C244-tyomnyy, RU-C244-svetlyy

\chapter{Ancient, Loud, Bright, Dark, Light}
\label{ch:ru-ancient-loud-bright-dark-light}
\hlchaptermodality{244}

\begin{chapteropening}
\textbf{By the end of this chapter:} \emph{I can say ancient, loud, bright, dark and light.}

Complete the last lesson of chapter 244: I can say ancient, loud, bright, dark and light.
\end{chapteropening}

\section[\texorpdfstring{\ru{древний} (drévniy)}{древний (drévniy)}]{\ru{древний} (drévniy) — ancient}

\label{lesson:RU-C244-drevniy}



\begin{quote}
Before the new one: say the Russian for low, then the Russian for ugly.
\end{quote}

\subsection*{You'll want to know: \ru{древний}}
\textbf{\ru{древний}} (\emph{drévniy}) — \textquotedblleft{}ancient\textquotedblright{}.

\begin{cousinweb}
\textbf{\ru{древний}} — Old Church Slavonic \emph{\ru{древьнь}}, \textquotedblleft{}ancient\textquotedblright{}, with the Church Slavonic \emph{-\ru{ре}-}. Vasmer derives it, with a \textquotedblleft{}perhaps\textquotedblright{}, from \emph{*drewos}, compare Gothic \emph{triggws}, \textquotedblleft{}faithful\textquotedblright{}, and Lithuanian \emph{drẽvė}, \textquotedblleft{}a hollow tree\textquotedblright{}.
\end{cousinweb}

\begin{grammarlens}[title={What you've built}]
One more word for this chapter.
\end{grammarlens}

\subsection*{Your turn}
\begin{itemize}
\raggedright
  \item \emph{Say it:} \emph{drévniy}
  \item \emph{Say it:} \emph{drévniy}, once more
  \item \emph{Say it:} say \emph{drévniy}
  \item \emph{Recall:} say the Russian for low, then the Russian for ugly, then say \emph{drévniy} again
\end{itemize}

\begin{tcolorbox}[breakable,colback=teal!4,colframe=teal!35!black,title={Before you move on}]
What does \ru{древний} mean? (\textquotedblleft{}Ancient\textquotedblright{}.) Say it once more.
\end{tcolorbox}
\section[\texorpdfstring{\ru{громкий} (grómkiy)}{громкий (grómkiy)}]{\ru{громкий} (grómkiy) — loud}

\label{lesson:RU-C244-gromkiy}



\begin{quote}
Before the new one: say the Russian for ugly, then the Russian for ancient.
\end{quote}

\subsection*{You'll want to know: \ru{громкий}}
\textbf{\ru{громкий}} (\emph{grómkiy}) — \textquotedblleft{}loud\textquotedblright{}.

\begin{cousinweb}
\textbf{\ru{громкий}} comes from \emph{\ru{гром}}, \textquotedblleft{}thunder\textquotedblright{} (Vasmer): loud as thunder.
\end{cousinweb}

\begin{grammarlens}[title={What you've built}]
One more word for this chapter.
\end{grammarlens}

\subsection*{Your turn}
\begin{itemize}
\raggedright
  \item \emph{Say it:} \emph{grómkiy}
  \item \emph{Say it:} \emph{grómkiy}, once more
  \item \emph{Say it:} say \emph{grómkiy}
  \item \emph{Recall:} say the Russian for ugly, then the Russian for ancient, then say \emph{grómkiy} again
\end{itemize}

\begin{tcolorbox}[breakable,colback=teal!4,colframe=teal!35!black,title={Before you move on}]
What does \ru{громкий} mean? (\textquotedblleft{}Loud\textquotedblright{}.) Say it once more.
\end{tcolorbox}
\section[\texorpdfstring{\ru{яркий} (yárkiy)}{яркий (yárkiy)}]{\ru{яркий} (yárkiy) — bright}

\label{lesson:RU-C244-yarkiy}



\begin{quote}
Before the new one: say the Russian for ancient, then the Russian for loud.
\end{quote}

\subsection*{You'll want to know: \ru{яркий}}
\textbf{\ru{яркий}} (\emph{yárkiy}) — \textquotedblleft{}bright\textquotedblright{}.

\begin{cousinweb}
\textbf{\ru{яркий}} goes with \emph{\ru{ярый}}, \textquotedblleft{}fierce, fiery\textquotedblright{} (Vasmer), Proto-Slavic \emph{*jarŭ}, which Vasmer relates to Greek \emph{zōrós}, \textquotedblleft{}fiery, strong\textquotedblright{}.
\end{cousinweb}

\begin{grammarlens}[title={What you've built}]
One more word for this chapter.
\end{grammarlens}

\subsection*{Your turn}
\begin{itemize}
\raggedright
  \item \emph{Say it:} \emph{yárkiy}
  \item \emph{Say it:} \emph{yárkiy}, once more
  \item \emph{Say it:} say \emph{yárkiy}
  \item \emph{Recall:} say the Russian for ancient, then the Russian for loud, then say \emph{yárkiy} again
\end{itemize}

\begin{tcolorbox}[breakable,colback=teal!4,colframe=teal!35!black,title={Before you move on}]
What does \ru{яркий} mean? (\textquotedblleft{}Bright\textquotedblright{}.) Say it once more.
\end{tcolorbox}
\section[\texorpdfstring{\ru{тёмный} (tyómnyy)}{тёмный (tyómnyy)}]{\ru{тёмный} (tyómnyy) — dark}

\label{lesson:RU-C244-tyomnyy}



\begin{quote}
Before the new one: say the Russian for loud, then the Russian for bright.
\end{quote}

\subsection*{You'll want to know: \ru{тёмный}}
\textbf{\ru{тёмный}} (\emph{tyómnyy}) — \textquotedblleft{}dark\textquotedblright{}.

\begin{cousinweb}
\textbf{\ru{тёмный}} is built on \emph{\ru{тьма}}, \textquotedblleft{}darkness\textquotedblright{}, Proto-Slavic \emph{*tĭma}, which Vasmer relates to Lithuanian \emph{tamsà}, Sanskrit \emph{támas}, \textquotedblleft{}darkness\textquotedblright{}, and Latin \emph{tenebrae}, \textquotedblleft{}darkness\textquotedblright{}.
\end{cousinweb}

\begin{grammarlens}[title={What you've built}]
One more word for this chapter.
\end{grammarlens}

\subsection*{Your turn}
\begin{itemize}
\raggedright
  \item \emph{Say it:} \emph{tyómnyy}
  \item \emph{Say it:} \emph{tyómnyy}, once more
  \item \emph{Say it:} say \emph{tyómnyy}
  \item \emph{Recall:} say the Russian for loud, then the Russian for bright, then say \emph{tyómnyy} again
\end{itemize}

\begin{tcolorbox}[breakable,colback=teal!4,colframe=teal!35!black,title={Before you move on}]
What does \ru{тёмный} mean? (\textquotedblleft{}Dark\textquotedblright{}.) Say it once more.
\end{tcolorbox}
\section[\texorpdfstring{\ru{светлый} (svétlyy)}{светлый (svétlyy)}]{\ru{светлый} (svétlyy) — light, bright (of a colour or room)}

\label{lesson:RU-C244-svetlyy}



\begin{quote}
Before the new one: say the Russian for bright, then the Russian for dark.
\end{quote}

\subsection*{You'll want to know: \ru{светлый}}
\textbf{\ru{светлый}} (\emph{svétlyy}) — \textquotedblleft{}light, bright (of a colour or room)\textquotedblright{}.

\begin{cousinweb}
\textbf{\ru{светлый}} $\leftarrow$ Proto-Slavic \emph{*světĭlŭ}, from \emph{\ru{свет}}, \textquotedblleft{}light\textquotedblright{} (Vasmer).
\end{cousinweb}

\begin{grammarlens}[title={What you've built}]
One more word for this chapter.
\end{grammarlens}

\subsection*{Your turn}
\begin{itemize}
\raggedright
  \item \emph{Say it:} \emph{svétlyy}
  \item \emph{Say it:} \emph{svétlyy}, once more
  \item \emph{Say it:} say \emph{svétlyy}
  \item \emph{Recall:} say the Russian for bright, then the Russian for dark, then say \emph{svétlyy} again
\end{itemize}

\begin{tcolorbox}[breakable,colback=teal!4,colframe=teal!35!black,title={Before you move on}]
What does \ru{светлый} mean? (\textquotedblleft{}Light, bright (of a colour or room)\textquotedblright{}.) Say it once more.
\end{tcolorbox}
